\scanpage{158}
% archive foliation: 157
$X$ schéma \uline{Steinien} sur $S$

\deleted{\unsure{Soit}} Soit $\uline L$ inversible\inserted{très ample} sur $X$, tel que $f_*(\uline L)\inserted{=E}$ [soit libre de rang $N$]\inserted{tel que $i:X-\mathrm P(\unsure{E})$}
\margin{\unsure{inutile?}}
\margin{\unsure{sens ordinaire}}
\margin{$A_1$}
Équivalence

(i) Donner une base de $E=f_*(\uline L)$ i.e. un \uline{isom.} $\uline O_S^N\simeq\unsure{E}$

(ii) Donner un isom. $\varphi:\mathrm P_S^N\xrightarrow{\sim}\mathrm P(E\deleted{\unsure{L\to E}})$, $\oplus$ un isom.
$(\varphi^{-1}i)^*(\uline O(1))\simeq\uline L$

(iii) Donner un \deleted{\unsure{immersion}} morphisme
\[
j:X-\mathrm P_S^N
\]
$\oplus$ un isomorphisme $\alpha:j^*(\uline O(1))\simeq\uline L$
tel que l'hom canonique
\[
q_*(\uline O(1))\longrightarrow f_*(\uline L)=E
\]
soit un isom. [condition qui ne dépend pas du choix de $\alpha$]

\margin{\unsure{Dans la définition de spécial, faire attention qu'on exige que le changement de base donne un} \uline{morphisme spécial}.
\uuline{N.} $\unsure{f'_*f'^*(\uline L)}\uuline{\unsure{\simeq}}\unsure{N'\otimes_A A'}$}
\margin{Compatibles aux \unsure{deux} \unsure{opérations de $\Pi(g^*)$} et localisation.}

Équivalences

(i) Donner une \deleted{base}\inserted{section} du faisceau des germes de bases de $E$ modulo $\uline O_S^*$

(ii) Donner un isom. \unsure{$1$}
\[
\varphi:\deleted{\unsure{X}}\mathrm P_S^N-\mathrm P(E)
\]
(iii) Donner un morphisme
\[
j:X\longrightarrow P
\]
tel que

$1^\circ)$ $j^*(\uline O(1))\underset S\simeq\uline L$ \uuline{relativement à $S$}

$2^\circ)$ L'hom. $\underbrace{q_*(\uline O(1))\xrightarrow{\sim}f_*(j^*(\uline O(1)))\ \text{un isom.}}_{\text{morphisme spécial}}$

